\chapter{Instruments}\label{ch:iv}
\newcommand{\ivwrightols}{$-$0.50}
\newcommand{\ivwrightiv}{$-$1.01}
\newcommand{\ivcardn}{3{,}010}
\newcommand{\ivcardols}{0.075}
\newcommand{\ivcardolsse}{0.003}
\newcommand{\ivcardfs}{0.32}
\newcommand{\ivcardfsse}{0.09}
\newcommand{\ivcardF}{13.3}
\newcommand{\ivcardtsls}{0.132}
\newcommand{\ivcardtslsse}{0.055}
\newcommand{\ivcardrfmed}{0.131}
\newcommand{\ivcardrflo}{0.043}
\newcommand{\ivcardrfhi}{0.249}
\newcommand{\ivcardrfpgt}{0.86}
\newcommand{\ivcardsmed}{0.135}
\newcommand{\ivcardslo}{0.039}
\newcommand{\ivcardshi}{0.270}
\newcommand{\ivcardspgt}{0.86}
\newcommand{\ivcardrho}{$-$0.29}
\newcommand{\ivcardrholo}{$-$0.71}
\newcommand{\ivcardrhohi}{0.17}
\newcommand{\ivweakF}{1.7}
\newcommand{\ivmanyF}{1.0}
\newcommand{\ivweakols}{0.70}
\newcommand{\ivweakiqrlo}{$-$0.43}
\newcommand{\ivweakiqrhi}{0.59}
\newcommand{\ivmanytsls}{0.66}
\newcommand{\ivweakbmed}{0.34}
\newcommand{\ivweakbcov}{98\%}
\newcommand{\ivmanybcov}{11\%}
\newcommand{\ivhiercov}{98\%}
\newcommand{\ivrflen}{7.3}
\newcommand{\ivrfcov}{98\%}
\newcommand{\ivrfmed}{0.29}
\newcommand{\ivweakblen}{5.2}
\newcommand{\ivmanybmed}{0.66}
\newcommand{\ivmanyblen}{0.60}
\newcommand{\ivhiermed}{0.58}
\newcommand{\ivhierlen}{3.1}
\newcommand{\ivhierrhat}{1.03}
\newcommand{\ivscalesmall}{14}
\newcommand{\ivscaleten}{1.9}
\newcommand{\ivscalehund}{0.77}
\newcommand{\ivscalehundcov}{4\%}
\newcommand{\ivcardrhopneg}{0.86}
\newcommand{\ivcardrhat}{1.000}
\newcommand{\ivcardess}{29{,}282}
\newcommand{\ivchexactmed}{0.12}
\newcommand{\ivchexactlen}{4.8}
\newcommand{\ivchexactcov}{100\%}
\newcommand{\ivchblockdiff}{0.01}
\newcommand{\ivcholdmed}{0.28}
\newcommand{\ivcholdlen}{0.78}
\newcommand{\ivcholdcov}{45\%}
\newcommand{\ivcholddiff}{0.22}
\newcommand{\ivcholdlonglen}{2.6}
\newcommand{\ivcholdlongcov}{80\%}
\newcommand{\ivchcardsd}{0.078}
\newcommand{\ivchcardoldrhat}{1.21}
\newcommand{\ivchcardoldess}{13}
\newcommand{\ivchcardoldlag}{0.995}
\newcommand{\ivchcardoldstep}{0.005}
\newcommand{\ivscaledeflen}{5.2}
\newcommand{\ivscaledefcov}{97\%}
\newcommand{\ivscalecovsmall}{97\%}
\newcommand{\ivscalecovone}{96\%}
\newcommand{\ivscalemedsmall}{0.20}
\newcommand{\ivscalemedhund}{0.64}
\newcommand{\ivhwtsls}{0.64}
\newcommand{\ivhwflatcov}{9\%}
\newcommand{\ivhwflatlen}{0.59}
\newcommand{\ivhwhiermed}{0.56}
\newcommand{\ivhwhiercov}{98\%}
\newcommand{\ivhwhierlen}{3.1}
\newcommand{\ivhwhierrhat}{1.03}
\newcommand{\ivhwhiersh}{99.7\%}
\newcommand{\ivhwhsmed}{0.59}
\newcommand{\ivhwhscov}{99\%}
\newcommand{\ivhwhslen}{5.5}
\newcommand{\ivhwhsrhat}{1.02}
\newcommand{\ivhwhssh}{99.8\%}
\newcommand{\ivhstsls}{0.17}
\newcommand{\ivhsflatmed}{0.17}
\newcommand{\ivhsflatcov}{77\%}
\newcommand{\ivhsflatlen}{0.25}
\newcommand{\ivhshiermed}{0.08}
\newcommand{\ivhshiercov}{88\%}
\newcommand{\ivhshierlen}{0.32}
\newcommand{\ivhshierrhat}{1.00}
\newcommand{\ivhshsmed}{0.10}
\newcommand{\ivhshscov}{89\%}
\newcommand{\ivhshslen}{0.30}
\newcommand{\ivhsrhat}{1.00}


\begin{tcolorbox}[colback=gray!6,colframe=gray!50,title=In brief,fonttitle=\normalfont\itshape,coltitle=black,colbacktitle=gray!15,boxrule=0.4pt,arc=1pt]
Instruments identify effects under unmeasured confounding at the price of an exclusion restriction. The local average effect is an accounting over compliers. With one weak instrument a posterior reports the weakness as a wide interval, whether the prior is on the reduced form or on the structural model, though the structural width is set by the prior on the error covariance. With twenty weak instruments a diffuse structural prior is confident and wrong, and a shrunk first stage gives intervals that cover the truth by being wide. Kernel methods are used for nonlinear structural functions, a global local prior is proposed for them, and Mendelian randomization turns Fisher's genetic confounder into an instrument.
\end{tcolorbox}

When confounders are unmeasured, adjustment fails\index{ignorability!failure of|textbf}, but another source of variation may rescue the analysis. An instrument\index{instrument|see{instrumental variable}}\index{instrumental variable|textbf} moves treatment, has no direct effect on the outcome, and is itself unconfounded. The idea is a century old.

\section{A short history of the instrument}\label{iv:sec:history}

\paragraph{Wright, 1928.} The instrumental variable\index{instrumental variable!history of|textbf} was invented for a question about tariffs. Appendix B of Philip Wright\index{Wright, Philip G.|textbf}'s \emph{The Tariff on Animal and Vegetable Oils} (Wright 1928) asks how to estimate supply and demand elasticities when price and quantity are determined together, so the observed points are intersections of two shifting curves and a regression of quantity on price traces neither. Wright\index{Wright, Philip G.!instrumental variable|textbf}'s answer was a variable that shifts one curve and not the other. A factor that moves supply, such as weather or yield per acre, while leaving demand alone moves the equilibrium along the demand curve, and the ratio of its covariances with quantity and with price recovers the demand slope. He applied the method to butter and flaxseed. Authorship has been debated, because Philip's son Sewall Wright\index{Wright, Sewall|textbf} had invented path analysis\index{path analysis} (Wright 1921, 1934); Stock and Trebbi (2003) attribute the appendix to Philip by stylometry, its path diagrams\index{Wright, Sewall!path diagrams|textbf} reflecting the collaboration of father and son.

Figure~\ref{fig:wright} reproduces Wright's problem. In 200 simulated markets demand is $q = 10 - p + u_d$ and supply is $q = 2 + 0.8p + 1.5w + u_s$, where the supply shifter $w$, call it weather, and the shocks $u_d$ and $u_s$ are independent standard normal. The cloud of equilibria slopes downward, but least squares gives \ivwrightols, a compromise between the two curves. The instrumental variable estimate, the covariance of weather with quantity divided by its covariance with price, is \ivwrightiv.

\begin{figure}[htbp]
\centering
\includegraphics[width=0.85\linewidth]{fig_wright.pdf}
\caption{Wright's problem. Equilibrium prices and quantities in 200 simulated markets, shaded by the supply shifter, darker for larger values. Least squares fits a compromise between the supply and demand curves; the instrumental variable estimate recovers the demand slope of $-1$.}
\label{fig:wright}
\end{figure}

\paragraph{Haavelmo and the Cowles Commission, 1943 to 1950.} The appendix was largely ignored, and the idea returned through simultaneous equations\index{simultaneous equations}. Haavelmo\index{Haavelmo, Trygve} (1943) showed that least squares in a simultaneous system is biased because regressor and error are correlated. Haavelmo (1944) treated econometric models as probability models, with autonomous structural equations and the reduced form into which they solve. In this framework exogenous variables play the role of Wright's shifters: one excluded from an equation identifies that equation's coefficients. The Cowles Commission\index{Cowles Commission|textbf} turned this into the theory of identification, with the order and rank conditions for linear systems and full information maximum likelihood (Koopmans, Rubin and Leipnik 1950), and limited information maximum likelihood (Anderson and Rubin 1949), whose accompanying test for the structural coefficient is still the weak instrument robust test of choice.

\paragraph{Reiers\o l, Geary and the name.} Olav Reiers\o l\index{Reiersol, Olav@Reiers\o l, Olav!the name instrumental variables|textbf} introduced the method in 1941 and the name instrumental variables\index{instrumental variable} in 1945, extending Frisch's (1934) confluence analysis of linear relations among mismeasured variables. Geary (1949) developed it for errors in variables\index{errors in variables|see{measurement error}}. The instrument correlates with the true regressor but not its measurement error.

\paragraph{Madansky.} Albert Madansky\index{Madansky, Albert!instruments for errors in variables|textbf} reviewed instruments alongside other methods for fitting lines when both variables have measurement error (Madansky 1959), then carried them to factor analysis (Madansky 1964), where indicators of one factor instrument another. For Madansky an instrument was an assumption about the covariance structure of unobservables, and its credibility a question about the science.

\paragraph{Two stage least squares and beyond.} Theil (1953) and Basmann (1957) introduced two stage least squares\index{two stage least squares|textbf}, and Sargan (1958) gave the general theory of instrumental variable estimation, including the test of overidentifying restrictions.  Imbens and Angrist (1994) and Angrist, Imbens and Rubin\index{Rubin, Donald B.!local average treatment effect|textbf} (1996) showed that with heterogeneous effects the instrument identifies the effect among compliers (Section~\ref{iv:sec:late}), and Bound, Jaeger and Baker (1995) and Staiger and Stock (1997) showed how badly the classical estimators behave when the instrument is weak.

\section{The linear model}
\begin{figure}[ht]
\centering
\begin{tikzpicture}[node distance=16mm]
\node[v] (Z) {$Z$};
\node[v,right=of Z] (T) {$T$};
\node[v,right=of T] (Y) {$Y$};
\node[h,above=10mm of $(T)!0.5!(Y)$] (U) {$U$};
\draw[e] (Z) -- (T); \draw[e] (T) -- (Y); \draw[e] (U) -- (T); \draw[e] (U) -- (Y);
\end{tikzpicture}
\caption{An instrument $Z$ moves treatment $T$, has no direct path to $Y$, and is independent of the unmeasured confounder $U$, shown dashed.}
\label{fig:dagiv}
\end{figure}


Consider the model of Figure~\ref{fig:dagiv},
\begin{align*}
Y &= \tau T + \delta U + \varepsilon, \\
T &= \pi Z + \gamma U + \nu,
\end{align*}
where $U$ is an unmeasured confounder and $Z$, $U$, $\varepsilon$ and $\nu$ are independent. Least squares of $Y$ on $T$ is biased by $\delta\gamma\,\mathrm{Var}(U)/\mathrm{Var}(T)$. The instrumental variable estimator\index{instrumental variable!linear model|textbf}
\[
\hat\tau_{\mathrm{IV}} = \frac{\widehat{\mathrm{Cov}}(Z, Y)}{\widehat{\mathrm{Cov}}(Z, T)}
\]
is consistent when $\pi \neq 0$, since $\mathrm{Cov}(Z, Y) = \tau\,\mathrm{Cov}(Z, T)$.

Figure~\ref{fig:iv} uses 500 replications of size 500 with $\tau = 1$, independent standard normal $Z$, $U$, $\nu$ and $\varepsilon$, $T = 0.5Z + U + \nu$ and $Y = \tau T + 2U + \varepsilon$. Here $\pi = 0.5$, $\gamma = 1$, $\delta = 2$ and $\mathrm{Var}(T) = 2.25$, so least squares has bias $2/2.25 = 0.89$ and averages 1.89. The instrumental variable estimate has median 1.01 and about four times the spread of least squares. With a first stage $F$ of about 60 it is roughly normal; Section~\ref{iv:sec:weak} shows the heavy tails of weak instruments.

\begin{figure}[ht]
\centering
\includegraphics[width=\linewidth]{iv.pdf}
\caption{Sampling distributions of least squares and instrumental variable estimates of $\tau = 1$ with an unmeasured confounder, 500 replications of size 500.}
\label{fig:iv}
\end{figure}

The IV\index{IV|see{instrumental variable}} estimator is a ratio, unstable when the denominator is near zero, and with one instrument and normal errors its finite sample distribution has no mean. 

\section{The local average treatment effect}\label{iv:sec:late}

With heterogeneous effects and binary $Z$ and $T$, let $T(z)$ be the potential treatment under instrument value $z$. Units are compliers\index{complier|textbf} if $T(1) = 1$ and $T(0) = 0$, always takers\index{always taker|textbf} if $T(1) = T(0) = 1$, never takers\index{never taker|textbf} if $T(1) = T(0) = 0$, and defiers\index{defier|textbf} if $T(1) = 0$ and $T(0) = 1$.

\begin{theorem}[Imbens and Angrist; Angrist, Imbens and Rubin]
Under random assignment of $Z$, the exclusion restriction\index{exclusion restriction|textbf} that $Z$ affects $Y$ only through $T$, monotonicity\index{monotonicity|textbf} (no defiers), and a nonzero first stage,
\[
\frac{\E(Y \mid Z = 1) - \E(Y \mid Z = 0)}{\E(T \mid Z = 1) - \E(T \mid Z = 0)} = \E\{Y(1) - Y(0) \mid \text{complier}\}.
\]
\end{theorem}

The left side is the Wald ratio\index{Wald estimator|textbf}\index{Wald ratio|see{Wald estimator}}; the right is the local average treatment effect\index{local average treatment effect|textbf}\index{LATE|see{local average treatment effect}}, the effect among those whose treatment the instrument moves. The intuition is an accounting. Random assignment makes the two arms comparable, so the numerator, the effect of $Z$ on $Y$, is a sum over the four strata. Always takers and never takers receive the same treatment under either value of the instrument, so by exclusion they contribute nothing; monotonicity removes the defiers, whose contribution would enter with the opposite sign. What remains is the complier share times the complier effect. The denominator, nonzero by assumption, is the complier share, the fraction of units the instrument moves. Dividing one by the other leaves the complier effect.

The local effect is an accounting identity\index{accounting identity} with $Q$ the covariate law of compliers, chosen by the instrument rather than the analyst, provided that, given covariates, compliers have the same average effect as everyone else. Report the complier share, because a Wald ratio built on a small share is a statement about few people, and the heavy tails of Figure~\ref{fig:weakiv} are near. Describe the compliers: when the assumptions hold given covariates $X$, their covariate law is identified as $Q_c(dx) \propto \pi_c(x)\,P(dx)$, where $\pi_c(x) = \E(T \mid Z = 1, X = x) - \E(T \mid Z = 0, X = x)$ is the complier share at $x$. And do not present the number as the average effect in the population unless a model extends it there (Section~\ref{iv:sec:accounting}).

\section{Bayesian instrumental variables}

Imbens and Rubin\index{Rubin, Donald B.!Bayesian instrumental variables|textbf} (1997) gave a Bayesian analysis of the binary instrument model\index{instrumental variable!Bayesian|textbf}. The principal strata, compliers, always takers and never takers, are latent, but $(Z, T)$ partly reveals them: a unit with $Z = 1$ and $T = 0$ must be a never taker, while one with $Z = 1$ and $T = 1$ is a complier or an always taker. Data augmentation imputes the missing strata, and the outcome model in each stratum is fitted to the completed data.

The Bayesian analysis respects the model's constraints: stratum proportions are nonnegative and outcome distributions proper, whereas the moment calculation behind the Wald ratio, applied to outcome distributions rather than means, can give complier densities that are negative in finite samples. It also delivers effects for always takers and never takers under explicit priors, separating what is learned from data and what from assumption.

\paragraph{Other Bayesian approaches.} The structural tradition models treatment and outcome jointly with correlated errors, so the posterior of the structural coefficient propagates first stage uncertainty rather than plugging in a fit. Kleibergen and Zivot (2003) compared diffuse priors in the structural and reduced form parametrizations, a comparison taken up in Section~\ref{iv:sec:weak}. Chib and Hamilton (2002) gave semiparametric Bayesian treatment models with an error distribution not fixed in advance. Conley, Hansen, McCulloch and Rossi (2008) replaced bivariate normal errors with a Dirichlet process mixture, avoiding a parametric error model. Conley, Hansen and Rossi (2012) placed a prior on the instrument's direct effect on the outcome, so exclusion need hold only approximately, the plausibly exogenous instruments of their title. The natural nonparametric extension puts BART or Gaussian process priors on both the first stage and the structural function, the joint posterior carrying uncertainty from one to the other.

\paragraph{Selection models.} The bivariate normal selection model of Chapter~\ref{ch:frameworks} (equation~\eqref{ss:eq:mills}) is the same structural idea with a latent selection equation in place of a linear first stage. Given the latent propensity $T^\ast$, the selection error $V$ is known, and the outcome equation is a regression with the added term $\rho\sigma V$ and residual variance $\sigma^2(1-\rho^2)$. Given the parameters and the outcome, each $T^\ast_i$ is normal with variance $1-\rho^2$ and mean equal to the selection index plus $\rho/\sigma$ times the outcome residual $Y_i - X_i^\top\beta - \tau T_i$, truncated to be positive when $T_i = 1$ and negative otherwise, so a Gibbs sampler alternates between the two (Chib and Hamilton 2002). With a weak or absent instrument the data barely identify the correlation $\rho$. Setting $\rho = 0$ is ignorability. Without an excluded instrument, what the data say about $\rho$ comes from the normal tails, and a prior on $\rho$ leaves that restriction in place. A prior concentrated near zero but not at it usually widens the posterior of $\tau$ when the data say little about $\rho$, but nothing guarantees it. Read this way, $\rho$ is a sensitivity parameter and the prior on it is the sensitivity analysis of Section~\ref{sec:sensitivity}: report how that posterior moves as the prior spreads.

\section{Weak instruments, Bayesian style}\label{iv:sec:weak}

When the first stage\index{first stage|textbf} is weak\index{weak instrument|textbf}, the instrumental variable\index{instrumental variable} estimator is badly biased toward least squares when there are many instruments and erratic when there is one, and in both cases its confidence intervals undercover. The Anderson and Rubin\index{Anderson and Rubin test|textbf} (1949) test of $\tau = \tau_0$ asks whether the instrument has any effect on $Y - \tau_0 T$ once the controls are in the regression. Its 90\% confidence set collects the values $\tau_0$ not rejected and has the stated coverage whatever the strength of the first stage, exactly with normal homoskedastic errors and approximately in large samples otherwise. With one instrument it is unbounded when the first stage $F$ statistic falls below the test's critical value. The Bayesian response\index{weak instrument!Bayesian treatment|textbf} is simpler. Write the reduced form\index{reduced form|textbf} as
\[
Y = \pi_Y Z + u_Y, \qquad T = \pi_T Z + u_T,
\]
so that $\tau = \pi_Y/\pi_T$. In the notation of the linear model $\pi_T = \pi$, $\pi_Y = \tau\pi$, $u_T = \gamma U + \nu$ and $u_Y = \tau u_T + \delta U + \varepsilon$. The posterior for $(\pi_Y, \pi_T)$ is well behaved whatever their size. The posterior for $\tau$ is the law of their ratio, wide and heavy tailed when $\pi_T$ is near zero, because the data then say little about $\tau$.

Kleibergen and Zivot showed that with weak instruments the parametrization matters: priors flat in the structural form can be informative in the reduced form, and the reverse. The reduced form is the natural place to start. Because $\pi_Y = \tau\pi_T$, a prior flat in $(\pi_Y, \pi_T)$ is proportional to $|\pi_T|$ in the coordinates $(\tau, \pi_T)$, so it puts little weight on a small first stage, whereas a prior flat in the structural coefficients weights that region fully, and near $\pi_T = 0$ the structural likelihood has a ridge along which the data cannot separate $\tau$ from the endogeneity (Section~\ref{iv:sec:many}).

Figure~\ref{fig:weakiv} shows a simulation in the spirit of quarter of birth. The true effect is $\tau = 0.1$, the treatment and outcome errors are standard normal with correlation $0.6$, there are 2{,}000 observations, the instruments are standard normal with first stage coefficient $0.03$ for the single instrument and $0.03/\sqrt{20}$ for each of the twenty (concentration parameter 1.8 in both, defined below), and they are very weak: one with median first stage $F$ of \ivweakF\ in the top panel, twenty with median $F$ of \ivmanyF\ in the bottom panel. Least squares centres on \ivweakols. With one weak instrument two stage least squares is centred near the truth but enormously variable, with interquartile range from \ivweakiqrlo\ to \ivweakiqrhi; with twenty it is concentrated near least squares, at \ivmanytsls, as Bound, Jaeger and Baker found.

 The structural model is that of Rossi, Allenby and McCulloch (2005): $T = \pi Z + W\gamma_T + e_1$ and $Y = \tau T + W\gamma_Y + e_2$, where $W$ holds the controls and the intercept and $(e_1, e_2)$ is bivariate normal with covariance $\Sigma$. The priors are $N(0, 10^4)$ on $\tau$, flat on the other coefficients, and inverse Wishart on $\Sigma$ with 3 degrees of freedom and a diagonal scale matrix holding one tenth of the sample variances of treatment and outcome. Integrating out everything but $\tau$ leaves a one dimensional posterior, which we evaluate on a grid, so these results carry no Monte Carlo error. The Gibbs samplers used below put $N(0, 10^4)$ priors, with standard deviation 100, on the other coefficients as well. We report the median across replications of the posterior median and of the length of the equal tailed 90\% interval, and the share of replications in which that interval covers the truth. With one weak instrument the structural posterior reports the weakness: the median of its posterior medians is \ivweakbmed, between the truth and least squares, its intervals have median length \ivweakblen, and they cover the truth in \ivweakbcov\ of replications. With a flat prior on the reduced form the posterior for $\tau = \pi_Y/\pi_T$ is wider still, with median interval length \ivrflen, and it covers the truth in \ivrfcov\ of replications. Both say that these data barely inform $\tau$.

How wide the structural posterior is depends on the prior on $\Sigma$. With a weak instrument the data fix little more than the least squares slope (Section~\ref{iv:sec:many} shows why), and along the ridge of pairs $(\tau, \Sigma)$ that fit the data about equally well the inverse Wishart prior acts approximately as a normal prior on $\tau$ centred at the least squares slope, with precision proportional to its scale (Exercise~\ref{ex:iv:iwscale}). On the first 100 data sets, scales of 0.01, 0.1, 1 and 10 times the sample variances give intervals of median length \ivscalesmall, \ivscaledeflen, \ivscaleten\ and \ivscalehund, which cover the truth in \ivscalecovsmall, \ivscaledefcov, \ivscalecovone\ and \ivscalehundcov\ of data sets. At the smallest scale the medians move toward the truth, to \ivscalemedsmall, inside very wide intervals; at the largest the posterior is confident and wrong, its medians near least squares at \ivscalemedhund. A prior on a nuisance covariance that looks harmless sets the width of the answer when the instrument is weak, and the analyst should show what happens when it changes.

With twenty weak instruments the diffuse structural posterior fails. Its medians sit at \ivmanybmed, next to two stage least squares at \ivmanytsls\ and least squares at \ivweakols, and its intervals have median length \ivmanyblen\ and cover the truth in \ivmanybcov\ of replications. A flat prior on twenty first stage coefficients is not uninformative about the strength of the instruments. The concentration parameter, the first stage signal $\pi^\top Z^\top Z\pi$ divided by the first stage error variance, is a quadratic form in the coefficients, and a flat prior on twenty of them gives it a density proportional to its ninth power (Exercise~\ref{ex:iv:flatK}): the prior says the first stage is strong, and the posterior behaves as if it were. With one instrument the same density is proportional to the inverse square root of the concentration parameter and leans toward a weak first stage. (The reduced form posteriors use the first 200 of the 300 data sets, the comparison of scales the first 100, and the other results all 300.) A posterior reports weakness only if its prior allows a weak first stage, as the shrinkage priors of Section~\ref{iv:sec:many} do.

\begin{figure}[htbp]
\centering
\includegraphics[width=\linewidth]{fig_weakiv.pdf}
\caption{Weak instruments in the spirit of quarter of birth: one instrument (top) and twenty (bottom), 300 data sets each. Gray: two stage least squares; black outline: posterior medians under the diffuse structural prior, computed exactly. The solid line is the truth $\tau = 0.1$, the dashed line the median least squares estimate. Estimates outside the plotted range are counted in the end bins.}
\label{fig:weakiv}
\end{figure}

\section{Many instruments and the horseshoe}\label{iv:sec:many}

In finite samples with many instruments\index{many instruments|textbf}, the fitted treatment inherits first stage error correlated with the structural error, even when the true first stage is zero, biasing two stage least squares\index{two stage least squares}\index{two stage least squares!bias toward least squares|textbf} toward least squares. Shrinking the first stage reduces this overfitting. Belloni, Chen, Chernozhukov and Hansen (2012) used the lasso; a horseshoe first stage\index{horseshoe prior!for many instruments|textbf}\index{horseshoe prior} leaves strong coefficients almost untouched and pulls weak ones to zero, avoiding the lasso's constant shrinkage of strong instruments. It is not a cure for weakness\index{weak instrument!and the horseshoe|textbf}. With a zero first stage, and approximately with a weak one, $T = u_T$. Write the structural outcome equation as $Y = \tau T + e$, where $e = u_Y - \tau u_T$ has standard deviation $\sigma_e$ and correlation $\rho$ with $u_T$, whose standard deviation is $\sigma_T$. Then $Y = (\tau + \rho\sigma_e/\sigma_T)\,T + w$ with $w$ uncorrelated with $T$ (independent under joint normality), so the data identify only the least squares slope, and the prior decides where $\tau$ sits along the ridge on which that slope is fixed. A prior that lets the first stage be near zero lets the posterior spread along the ridge; one that insists on a strong first stage, as a flat prior on twenty coefficients does, holds it near two stage least squares.

Table~\ref{tab:ivmany} compares two stage least squares and three first stage priors on six designs with twenty instruments, 1,000 data sets each, otherwise as in Figure~\ref{fig:weakiv}. The concentration parameter $c$ is 1.8, 20 or 135, carried by $s = 3$ of the instruments or spread over all twenty, each active instrument having first stage coefficient $\sqrt{c/(ns)}$. The weak design of that figure is $c = 1.8$ over twenty instruments, and the sparse design discussed below has $c = 135$ on three instruments with coefficient 0.15 and seventeen at zero. The hierarchical normal prior, in the spirit of Chamberlain and Imbens (2004), gives the twenty coefficients a common normal distribution whose variance has an inverse gamma prior with shape 1 and scale $10^{-4}$; the horseshoe gives each coefficient its own half Cauchy local scale under a half Cauchy global scale $\gscale$, both with unit scale. Neither prior is scaled by the error variance, and both keep the $N(0, 10^4)$ priors on $\tau$ and the controls and the prior on $\Sigma$ of Section~\ref{iv:sec:weak}. Both are fitted by the blocked Gibbs sampler described below, four chains of 4{,}000 iterations with the first 1{,}000 discarded, rerun at 16{,}000 with 4{,}000 discarded when the split R-hat of $\tau$ exceeds 1.01, and the diffuse prior is computed exactly as before.

\begin{table}[ht]
\centering
\caption{Twenty instruments, 1,000 data sets per design, $n = 2{,}000$, $\tau = 0.1$ and error correlation 0.6. The concentration parameter $c$ is carried by $s$ active instruments, each with first stage coefficient $\sqrt{c/(ns)}$; the weak design of the text is $c = 1.8$ with $s = 20$, the sparse design $c = 135$ with $s = 3$. For two stage least squares the median estimate over data sets; for each first stage prior the median over data sets of the posterior median, the share of equal tailed 90\% intervals that cover the truth, and the median interval length. The Monte Carlo standard error of a coverage near 90\% is about one point.}
\label{tab:ivmany}
\small
\setlength{\tabcolsep}{4pt}
\begin{tabular}{lcccccc}
\toprule
& \multicolumn{3}{c}{Three active, $s = 3$} & \multicolumn{3}{c}{Twenty active, $s = 20$} \\
\cmidrule(lr){2-4}\cmidrule(lr){5-7}
First stage prior & Median & Coverage & Length & Median & Coverage & Length \\
\midrule
\multicolumn{7}{@{}l}{\textit{Concentration $c = 1.8$}} \\
Two stage least squares & 0.66 &  &  & 0.64 &  &  \\
Diffuse, exact & 0.66 & 10\% & 0.59 & 0.64 & 9\% & 0.59 \\
Hierarchical normal & 0.59 & 98\% & 3.09 & 0.56 & 98\% & 3.06 \\
Horseshoe & 0.62 & 99\% & 5.54 & 0.59 & 99\% & 5.54 \\
\midrule
\multicolumn{7}{@{}l}{\textit{Concentration $c = 20$}} \\
Two stage least squares & 0.39 &  &  & 0.40 &  &  \\
Diffuse, exact & 0.37 & 40\% & 0.45 & 0.38 & 39\% & 0.45 \\
Hierarchical normal & $-0.07$ & 91\% & 1.77 & $-0.05$ & 91\% & 1.75 \\
Horseshoe & 0.08 & 92\% & 1.23 & 0.11 & 95\% & 1.76 \\
\midrule
\multicolumn{7}{@{}l}{\textit{Concentration $c = 135$}} \\
Two stage least squares & 0.17 &  &  & 0.17 &  &  \\
Diffuse, exact & 0.17 & 77\% & 0.25 & 0.17 & 76\% & 0.25 \\
Hierarchical normal & 0.08 & 88\% & 0.32 & 0.08 & 89\% & 0.32 \\
Horseshoe & 0.10 & 89\% & 0.30 & 0.10 & 90\% & 0.31 \\
\bottomrule
\end{tabular}
\end{table}


In the weak design the diffuse prior is again confident and wrong. Both shrinkage priors gain coverage by widening the intervals, with medians little closer to the truth than two stage least squares: the posterior has learned that the instruments are weak, not where $\tau$ is. In the sparse design the first stage is strong in total, yet two stage least squares and the diffuse posterior are still pulled toward least squares. The hierarchical normal prior and the horseshoe bring the medians toward the truth and improve the coverage, with intervals a little longer. This is the design in which the horseshoe's ability to keep a few strong coefficients while zeroing the rest should matter most, and it gains modestly for $\tau$: its posterior median has root mean squared error 0.094 against 0.105 for the hierarchical normal prior and mean error $-0.010$ against $-0.027$. It also estimates the first stage better, with half the error on the seventeen inactive instruments (root mean squared error 0.008 against 0.016). The larger gain comes at $c = 20$. With three and with twenty active instruments its posterior median has mean error $-0.07$ and $-0.04$ against $-0.22$ and $-0.19$ for the hierarchical prior, root mean squared error 0.34 and 0.37 against 0.46 and 0.46, and its intervals cover 0.922 and 0.945 against 0.908 and 0.906, with three active instruments on a shorter interval, 1.23 against 1.77. The samplers mixed: in the weak design the largest split R-hat over data sets was \ivhwhierrhat\ for the hierarchical prior and \ivhwhsrhat\ for the horseshoe, with \ivhwhiersh\ and \ivhwhssh\ of runs below 1.01 after the longer rerun that about 3 percent of them needed, and in the sparse design \ivhshierrhat\ and \ivhsrhat. At $c = 20$ the largest values reached 1.17, 5 to 8 percent of runs needed the longer chains, and at most 0.6 percent stayed above 1.01.

The structural model treats both stages jointly, its bivariate error $(e_1, e_2)$ being $(u_T, e)$ above, with a correlation $\rho$ that measures the endogeneity. Gibbs sampling for it is easy to write and easy to get wrong. The natural sampler alternates the structural coefficients, first stage coefficients and $\Sigma$, each conditional on the others. On the ridge each block is pinned by the others, and the sampler crawls. On the Card data of Section~\ref{iv:sec:card}, four chains of 6{,}000 iterations from different starting values gave a split R-hat of \ivchcardoldrhat\ for $\tau$ and an effective sample size of \ivchcardoldess\ from the 20{,}000 draws kept, with lag one autocorrelation \ivchcardoldlag; a typical step moved $\tau$ by \ivchcardoldstep, against a posterior standard deviation of \ivchcardsd. On the first 20 data sets of the one instrument design, runs of 1{,}200 iterations gave intervals of median length \ivcholdlen\ that covered the truth in \ivcholdcov\ of data sets, where the exact posterior gives \ivchexactlen\ and \ivchexactcov, and runs of 20{,}000 iterations still gave only \ivcholdlonglen\ and \ivcholdlongcov. Drawing $\tau$ in one block with the coefficient of the regression of $e_2$ on $e_1$ lets the sampler move along the ridge: on the same data sets its posterior medians differ from the exact ones by \ivchblockdiff\ at the median. The samplers behind Table~\ref{tab:ivmany} and the right panel of Figure~\ref{fig:card} are blocked, with the shrinkage scales of the first stage updated by their own conditional step. The posterior for the error correlation reports how much endogeneity the data see, and an analyst with many candidate instruments should report which ones the posterior keeps, since a design that rests on one or two of them is a different design from one that rests on twenty.

\section[Worked example: college proximity]{Worked example: college proximity and the return to schooling}\label{iv:sec:card}

Card\index{Card, David|see{college proximity instrument}} (1995) asked whether growing up near a four year college raises schooling and, through it, earnings. The data are the 1976 wave of the National Longitudinal Survey of Young Men, \ivcardn\ men with wages, schooling and an indicator of a four year college in the local labour market in 1966. We follow the textbook specification of Wooldridge (2020, Example 15.4), which treats experience as exogenous; Card also instruments experience with age. The instrument is proximity\index{college proximity instrument|textbf}; the controls are experience and its square, race, residence in a metropolitan area and in the South in 1976, metropolitan residence in 1966, and region of residence in 1966. The motivating worry is ability bias: men with more schooling may have earned more anyway.

Least squares gives a return of \ivcardols\ per year of schooling\index{college proximity instrument!return to schooling|textbf} (standard error \ivcardolsse). The first stage coefficient of proximity on schooling is \ivcardfs\ years (standard error \ivcardfsse), a first stage $F$ statistic of \ivcardF, just above the conventional threshold of 10. Two stage least squares gives \ivcardtsls\ (standard error \ivcardtslsse), larger than least squares, which surprised many: if ability bias pushes least squares up, the instrument should pull the estimate down. The usual explanations are unusually high returns among the compliers, men kept in school by a nearby college, and measurement error\index{measurement error|textbf} in schooling, which biases least squares toward zero. Schooling takes many values, so the instrumental variable estimate is not the binary complier effect of the theorem but the average causal response\index{average causal response|textbf} of Angrist and Imbens (1995), a weighted average of the effects of an extra year of schooling at each margin the instrument moves, and the complier share of Section~\ref{iv:sec:late} has no direct analogue.

Two Bayesian analyses show how much the parametrization matters at this strength of instrument. The first places flat priors on the coefficients of proximity and the controls in the reduced form wage and schooling equations and the prior $|\Omega|^{-3/2}$ on the covariance $\Omega$ of their errors, and takes the return as the ratio of the two proximity coefficients: posterior median \ivcardrfmed, 90\% interval from \ivcardrflo\ to \ivcardrfhi, and posterior probability \ivcardrfpgt\ that the return exceeds least squares. That interval agrees, to within simulation error, with the 90\% Anderson and Rubin confidence set of Section~\ref{iv:sec:weak}, which runs from 0.044 to 0.249. The second is the structural posterior of Section~\ref{iv:sec:weak}, computed exactly: median \ivcardsmed, 90\% interval from \ivcardslo\ to \ivcardshi, and probability \ivcardspgt\ of exceeding least squares. Figure~\ref{fig:card} shows it. The two agree on the centre and on the comparison with least squares, and the structural posterior has the longer tails. With a first stage $F$ of \ivcardF\ the parametrization matters little; it is with a weak instrument, as in Figure~\ref{fig:weakiv}, that the prior takes over.

\begin{figure}[htbp]
\centering
\includegraphics[width=\linewidth]{fig_card.pdf}
\caption{College proximity as an instrument for schooling (Card 1995). Left: exact posterior of the return to a year of schooling under the diffuse structural prior of Section~\ref{iv:sec:weak}, with the least squares (dashed) and two stage least squares (dotted) estimates. Right: posterior of the correlation between the errors of the schooling and wage equations, from the blocked Gibbs sampler of Section~\ref{iv:sec:many} (four chains of 11{,}000 iterations, the first 1{,}000 of each discarded; split R-hat \ivcardrhat\ and effective sample size \ivcardess\ for the return).}
\label{fig:card}
\end{figure}

The posterior of the error correlation has mean \ivcardrho, 90\% interval from \ivcardrholo\ to \ivcardrhohi, and probability \ivcardrhopneg\ of being negative. A negative correlation makes least squares too small, the direction that measurement error in schooling would produce, and the data lean that way, but they leave real probability on the positive correlation that ability bias implies. That is what comparing least squares with two stage least squares was trying to say, now with a probability attached.

\section{Instruments in unlikely places}\label{iv:sec:unlikely}

The search for instruments\index{instrumental variable!in unlikely places|textbf} has produced many ingenious and contested designs. A good instrument must move treatment for reasons unrelated to the outcome, so it is usually something nobody would think of as a cause of the outcome. Table~\ref{tab:unlikely} lists examples; most have drawn a critique of the assumption that made them clever.

\begin{table}[htbp]
\centering\scriptsize
\renewcommand{\arraystretch}{1.1}
\begin{tabular}{@{}L{1.0in}L{0.9in}L{0.95in}L{1.15in}@{}}
\toprule
Instrument & Treatment & Outcome & The worry \\
\midrule
Quarter of birth (Angrist and Krueger 1991) & years of schooling & earnings & weak first stage (Bound, Jaeger and Baker 1995); season of birth and family background (Buckles and Hungerman 2013) \\
Draft lottery number (Angrist 1990) & military service in Vietnam & civilian earnings & the lottery may change behaviour directly, through schooling to avoid service \\
Two first children of the same sex (Angrist and Evans 1998) & a third child & mother's labour supply & sex composition may affect costs directly \\
College nearby (Card 1995) & schooling & earnings & families choose where to live \\
Stormy weather at sea (Angrist, Graddy and Imbens 2000) & price of fish & quantity demanded & Wright's supply shifter again; few observations \\
Rainfall (Miguel, Satyanath and Sergenti 2004) & economic growth & civil conflict & rain may affect conflict directly (Sarsons 2015) \\
Settler mortality (Acemoglu, Johnson and Robinson 2001) & institutions & income per head & measurement of mortality (Albouy 2012) \\
Electoral cycles (Levitt 1997) & police numbers & crime & not significant once a coding error is corrected (McCrary 2002) \\
Streams that divide school districts (Hoxby 2000) & school choice & achievement & counting of streams (Rothstein 2007) \\
Randomly assigned judges (Kling 2006) & incarceration & later earnings & monotonicity across judges (Frandsen, Lefgren and Leslie 2023) \\
Genetic variants (Davey Smith and Ebrahim 2003) & an exposure such as alcohol & disease & pleiotropy (Section~\ref{iv:sec:mr}) \\
\bottomrule
\end{tabular}
\caption{Instruments in unlikely places. In each case the instrument is plausible precisely because it seems unrelated to the outcome, and most of the critiques question that.}
\label{tab:unlikely}
\end{table}



\paragraph{Quarter of birth.} Angrist and Krueger (1991) noticed that compulsory schooling laws require attendance until a birthday, usually the sixteenth, while school entry is by calendar year. Children born early in the year reach the leaving age with slightly less schooling than those born late, so quarter of birth\index{quarter of birth instrument|textbf} shifts schooling a little. Nobody believes season of birth affects earnings directly, which makes the instrument persuasive. But the first stage is tiny, and with many interactions of quarter of birth with state and year the instruments are individually very weak. Bound, Jaeger and Baker (1995) showed that two stage least squares with such instruments is biased toward least squares, and that randomly generated quarters of birth produced estimates and standard errors similar to the published ones.  Later work related season of birth to mothers' characteristics (Buckles and Hungerman 2013), so exclusion is not as safe as it looked either.

\paragraph{Same sex siblings.} Angrist and Evans (1998) used the fact that American parents whose first two children share a sex are more likely to have a third. A child's sex is as good as random, and it seems absurd that two daughters rather than a daughter and a son would affect a mother's employment except through family size. The design gives a credible local effect of a third child for families whose fertility responds to sex composition, a specific and perhaps unusual group of compliers.

\paragraph{Fish and weather.} Angrist, Graddy and Imbens (2000) used stormy weather to instrument the price of whiting in New York's Fulton fish market (Graddy 1995), repeating Wright's supply shifter: weather shifts supply, plausibly not demand. With only a few months of trading days, the design calls for uncertainty statements that hold in small samples.

\section{Nonparametric instruments}\label{iv:sec:nonpar}

With continuous treatment and a nonlinear structural function\index{nonparametric instrumental variables|textbf},
\[
Y = f(T) + \varepsilon, \qquad \E(\varepsilon \mid Z) = 0,
\]
and $f$ solves the integral equation
\[
\E(Y \mid Z) = \E\{f(T) \mid Z\} = (\mathcal{K} f)(Z),
\]
where $\mathcal{K}$ is the conditional expectation operator from functions of $T$ to functions of $Z$. The equation identifies $f$ only if $\mathcal{K}$ is injective, the completeness condition of Newey and Powell (2003): the only function $h$ with $\E\{h(T) \mid Z\} = 0$ is $h = 0$. Even then inverting $\mathcal{K}$ is ill posed\index{ill posed inverse problem|textbf}: $\mathcal{K}$ is typically compact, its singular values accumulate at zero, and small errors in $\E(Y \mid Z)$ produce large errors in $f$.

Kernel instrumental variables\index{kernel instrumental variable|textbf}, introduced by Singh, Sahani and Gretton in 2019, represent $f$ in a reproducing kernel Hilbert space and estimate $\mathcal{K}$ by kernel ridge regression of features of $T$ on features of $Z$; the second stage regresses $Y$ on the estimated conditional mean embedding. The first stage is a vector valued regression and the second a scalar one, with theory resting on source conditions that measure the smoothness of $f$ relative to $\mathcal{K}$.

\paragraph{Source conditions, priors and saturation.}  The H\"older source condition of order $r$, that $f = \mathcal{C}^{r} g$ with $\|g\| < \infty$ for the covariance operator $\mathcal{C}$, says that $f$ lies in the Cameron Martin space of the Gaussian prior with covariance $\mathcal{C}^{2r}$ (a prior that exists when $\mathcal{C}^{2r}$ has finite trace): in the eigenbasis of $\mathcal{C}$, with eigenvalues $s_k$, the Cameron Martin norm $\sum_k f_k^2/s_k^{2r}$ is $\|g\|^2$ (Exercise~\ref{ex:ch10:5}). When $\mathcal{C}$ has infinite rank, draws from that prior fall outside the space almost surely, since their norm is a sum of infinitely many squared independent standard normals. With fixed scale $c$ the posterior still contracts, but more slowly; attaining the rate below requires $c$ to depend on $n$. In the sequence model $y_k = f_k + n^{-1/2}\xi_k$, with $\xi_k$ independent standard normal, the consequences take a few lines:
\begin{itemize}
\item Ridge regression with penalty $\alpha$ multiplies $y_k$ by $s_k/(s_k + \alpha)$. That is the posterior mean under $f_k \sim N(0, cs_k)$ with $\alpha = 1/(nc)$, so the prior implicit in ridge has covariance proportional to $\mathcal{C}$, the case $r = 1/2$ of the reading above.
\item The prior with covariance $c\,\mathcal{C}^{2r}$ multiplies $y_k$ by $s_k^{2r}/(s_k^{2r} + a)$ with $a = 1/(nc)$. Under the source condition of order $r$, its squared bias is at most $a\|g\|^2/4$, since $as_k^{r}/(s_k^{2r} + a) \le \sqrt{a}/2$. With $s_k$ of order $k^{-b}$ and $2rb > 1$, the variance of the posterior mean and the posterior spread are both of order $a^{-1/(2rb)}/n$. Balancing gives posterior expected squared error, averaged over the data, of order $n^{-2rb/(2rb+1)}$, with $a$ of that order and $c$ of order $n^{-1/(2rb+1)}$. The contraction rate in norm is its square root, $\varepsilon_n = n^{-rb/(2rb+1)}$: writing $f_0$ for the truth, Markov's inequality gives $\E_{f_0}\Pi(\|f - f_0\| > M_n\varepsilon_n \mid y) \to 0$ for every $M_n \to \infty$, uniformly over $\|g\| \le R$ for fixed $R$. For $r \le 1$ ridge reaches the same order of worst case mean squared error in this model at the best $\alpha$ (squared bias of order $\alpha^{2r}$ from Exercise~\ref{ex:ch10:6}, variance of order $\alpha^{-1/b}/n$). With random design, taking $\mathcal{C}$ to be the kernel integral operator on $L^2$ of the law of $T$, Theorem 1 of Caponnetto and De Vito (2007) gives the same power for the squared prediction error of kernel ridge when $1/2 < r \le 1$ and $b > 1$, under a bounded kernel and moment conditions on the noise; at $r = 1/2$ their upper bound is $(\log n/n)^{b/(b+1)}$, larger by a logarithmic factor. 
\item Saturation, ridge regression ceasing to improve once the source condition exceeds a threshold, belongs to the fixed prior. With covariance proportional to $\mathcal{C}^{2r_0}$ and truth of order $r$, the same calculation gives squared bias of order $a^{2\min\{r/(2r_0),\,1\}}$ while the variance does not depend on $r$, so the rate stops improving at $r = 2r_0$. Ridge has $r_0 = 1/2$ and stops at $r = 1$ (Exercise~\ref{ex:ch10:6}). A Gaussian prior matched to the smoothness, $r_0 = r$, is below its ceiling for every $r$, but it has to be told $r$.
\item For unknown $r$ we propose a global local prior on a finite expansion: $f = \sum_{k=1}^m f_k\phi_k$, with $\phi_k$ the eigenfunctions of $\mathcal{C}$ and horseshoe scales on the coefficients, $f_k \mid \lambda_k, \gscale \sim N(0, \lambda_k^2\gscale^2)$, as in Chapter~\ref{ch:horseshoe}. The cutoff $m$ is part of the prior: infinitely many coefficients that are independent and identically distributed given $\gscale$ have $\sum_k f_k^2 = \infty$ almost surely, so without it the prior does not live on the space of the source condition. With $m$ growing in $n$, and eigenvalues in decreasing order, a truth of order $r$ has squared truncation error at most $s_{m+1}^{2r}\|g\|^2$. We expect the prior to adapt without saturating, and we have no proof and no rule for choosing $m$. The expectation rests on where the ceiling comes from: in Exercise~\ref{ex:ch10:6} it is the bias on the largest coordinates, of order $\alpha$ for ridge, while the horseshoe's tails shrink a large $y_k$ by only about $2/(ny_k)$ (Chapter~\ref{ch:horseshoe}). The contraction results we know for the horseshoe, such as van der Pas, Szab\'o and van der Vaart (2017a), concern nearly black vectors, not source sets, so whether it attains the rate for every $r$ remains open.
\end{itemize}
With an operator in the model the coordinates are observed through its singular values, which slow the rate. Knapik, van der Vaart and van Zanten (2011) carry out the calculation for Gaussian priors: with $\mathcal{K}$ known, diagonal in the eigenbasis of $\mathcal{C}$, with singular values of order $k^{-p}$, and with $c$ of order $n^{-1/(2rb+2p+1)}$, the rate in norm is $n^{-rb/(2rb+2p+1)}$, and $p = 0$ returns the rate above. This is a benchmark, since in nonparametric instruments $\mathcal{K}$ is estimated and its error needs control of its own. A rate is also not coverage: the same paper shows that credible sets under a prior smoother than the truth can have frequentist coverage tending to zero. Knapik, Szab\'o, van der Vaart and van Zanten (2016) show that empirical and hierarchical Bayes over the regularity of a Gaussian prior adapt in mildly ill posed problems. The horseshoe would have to do at least as well to justify its place.

\paragraph{Scalability.} Kernel methods scale as the cube of the sample size, prohibitive for instruments in large administrative data sets. Random features, Nystr\"om approximation and conjugate gradient solvers cut the cost, and the first two are themselves priors. With $m$ random features $\varphi_j$, the function $f = \sum_j \beta_j\varphi_j$ with $\beta_j$ independent $N(0, 1/m)$ has a Gaussian prior of dimension $m$ and covariance $m^{-1}\sum_j \varphi_j(t)\varphi_j(t')$. Nystr\"om at $m$ points $v$ replaces the kernel $\kappa$ by $\kappa(t, v)\kappa(v, v)^{-1}\kappa(v, t')$, the covariance of $\E\{f(t) \mid f(v)\}$ under the Gaussian process prior with kernel $\kappa$, a prior covariance of rank $m$. Either way the intervals are those of the low rank prior, and they should not be read as those of the full kernel. The global local prior is conditionally Gaussian: given its scales, an EM step or a Gibbs sweep is a weighted ridge regression, at about the cost of one kernel ridge regression. How many such steps a fit takes is an open question for us.

\section{Mendelian randomization}\label{iv:sec:mr}

Within families, genetic variants are assigned at conception almost randomly; those that affect an exposure can serve as instruments for it. This is Mendelian randomization\index{Mendelian randomization|textbf}, proposed by Katan (1986) and developed by Davey Smith and Ebrahim (2003).

Fisher's constitutional hypothesis (Chapter~\ref{ch:smoking}) was that genes caused both smoking and lung cancer. Mendelian randomization turns it around: a gene that affects smoking quantity, and lung cancer only through smoking, is an instrument for smoking. Variants in the nicotinic receptor gene cluster on chromosome 15, which are associated with the number of cigarettes smoked per day (Thorgeirsson et al. 2008), have been used in this way to study the effects of smoking heaviness.

Exclusion is the weak point. A variant may affect the outcome through pathways other than the exposure, which geneticists call horizontal pleiotropy\index{Mendelian randomization!pleiotropy|textbf}. With many variants, MR Egger regression (Bowden, Davey Smith and Burgess 2015), which needs the direct effects to be independent of the variants' effects on the exposure, and weighted median estimators (Bowden, Davey Smith, Haycock and Burgess 2016), which need more than half of the weight to come from valid variants, relax the requirement that every variant be valid. A Bayesian treatment puts a sparsity prior on variants' direct effects, most near zero and a few not. Provided most variants are valid, a horseshoe\index{horseshoe prior} lets that majority locate the effect; variants whose direct effects escape shrinkage look invalid. The horseshoe gives no exact zeros, so a posterior probability that a variant is valid needs a prior with a point mass at zero.

\section{Instruments as accounting}\label{iv:sec:accounting}

Two valid instruments\index{instrumental variable!accounting view|textbf} for the same treatment can give different local average treatment effects\index{local average treatment effect}, both correct, because each defines its own compliers. On the accounting view\index{accounting identity!and instruments|textbf} the object of inference is the structural model: in the binary case, the joint law of potential treatments and outcomes. The local effect\index{local average treatment effect!as accounting|textbf} is one accounting of it, over that instrument's compliers. A model of how effects and compliance vary with covariates can compute the local effect for any instrument and the average effect for any population. Without one, each instrument answers only its own question.

College proximity moves men who would otherwise have stopped at the end of secondary school; quarter of birth moves those who leave at the earliest legal age. Both instrument schooling, yet their compliers can have different returns even if both exclusion restrictions hold. A joint model of returns and compliance computes each local effect by averaging the same conditional effect with different weights, provided compliance carries no information about the effect once covariates are fixed. Population weights give the population average. Where compliers cover only part of the covariate space, this average relies on extrapolation; posterior width there comes from the prior on how effects vary with covariates, so it should be reported under more than one such prior.

Report the local effect and its compliers, and, when a model is fitted, the average for the population of interest with its wider interval, showing how much comes from data and how much from the model.

\section{In practice}

\paragraph{The target.} Each posterior draw gives $\tau^{(s)}$, directly or as $\pi_Y^{(s)}/\pi_T^{(s)}$; the complier effect from the outcome models of the imputed strata; or $f^{(s)}$ from the eigenbasis coefficients. Averaging $f^{(s)}(t) - f^{(s)}(t')$ over any law of $(t, t')$ gives a policy effect, again by accounting.

\paragraph{What the analyst does differently.} An analyst with an instrument should state the exclusion restriction as a scientific claim and say who the compliers are, since the local effect describes no one else. She should report the first stage with its posterior, not only an $F$ statistic, and the posterior for the error correlation, since it reports the endogeneity the design exists to address. With many instruments she should shrink the first stage and report which instruments the posterior keeps (Section~\ref{iv:sec:many}). She should state the prior on the error covariance and report the posterior under more than one prior, because with one weak instrument the scale of that prior moved the coverage of nominal 90\% intervals from \ivscaledefcov\ to \ivscalehundcov\ (Section~\ref{iv:sec:weak}). She should also run more than one chain from different starting values, because a one at a time Gibbs sampler crawls on the ridge (Section~\ref{iv:sec:many}).

\paragraph{Against DML.} Double machine learning has an instrumental variable score (Chernozhukov et al. 2018; Table~\ref{tab:dml:scores}), and for a single strong instrument in a linear model it and two stage least squares are adequate and cheaper; the Bayesian answer then agrees with two stage least squares, as in the Card data, and costs more, so the case for it lies elsewhere. We did not run the score, on the Card data or on the designs of Section~\ref{iv:sec:weak}, so what the runs show concerns two stage least squares and the posteriors. With one weak instrument two stage least squares is erratic, while the structural and reduced form posteriors report the weakness as wide intervals. With twenty weak instruments no method we ran recovers the truth: two stage least squares and the diffuse structural posterior sit next to least squares, and the shrunk first stages of Table~\ref{tab:ivmany} cover the truth by being wide, not by being right. With a nonlinear structural function the ill posed inverse needs a prior matched to the unknown smoothness, and the global local prior of Section~\ref{iv:sec:nonpar} is proposed for that purpose; we have not tested it here. The limit on our side is the weight a weak instrument puts on the prior for the error covariance: the Bayesian analyst must state that prior and show what changes when it does.

\section*{Notes}



The chapter assumes monotonicity; Balke and Pearl (1997) bound the average effect without it, by linear programming over response types, as in Chapter~\ref{ch:tianpearl}.

\begin{exercises}
\item\label{ex:ch10:1} Derive the bias of least squares in the linear model and check it against the simulated mean of 1.89 in the design of Figure~\ref{fig:iv}.
\item\label{ex:ch10:2} Show that the Wald ratio with one instrument has no finite mean when the reduced form and first stage coefficient estimates are jointly normal, the first stage estimate has positive variance and the correlation between the two estimates is below one in absolute value.
\item\label{ex:ch10:3} Prove the theorem of Section~\ref{iv:sec:late}, and show where each assumption is used.
\item\label{ex:ch10:4} Write the data augmentation sampler for the Imbens and Rubin model with normal outcomes in each stratum.
\item\label{ex:ch10:5} Show that the H\"older source condition $f = \mathcal{C}^{r}g$ with $\|g\| < \infty$ says that $f$ lies in the Cameron Martin space of the Gaussian prior with covariance $\mathcal{C}^{2r}$, assuming $\mathcal{C}^{2r}$ has finite trace.
\item\label{ex:ch10:6} Explain why ridge regression saturates at source condition $r = 1$.
\item\label{ex:iv:wright} In Wright's market with demand $q = a - bp + u_d$ and supply $q = c + dp + ew + u_s$, derive the probability limit of the least squares slope of $q$ on $p$ and show that it is a weighted average of $-b$ and $d$. Check the value against Figure~\ref{fig:wright}, which shows a single sample of 200 markets.
\item\label{ex:iv:card} Using \texttt{iv\_card.py} and \texttt{iv\_weak2.py}, add the father's and mother's education to the controls in the Card example, which loses observations. How do two stage least squares, the reduced form posterior and the exact structural posterior change, and why?
\item\label{ex:iv:qob} Explain why replacing true quarters of birth by random numbers would, with many interactions, give two stage least squares estimates close to least squares. What does the reduced form posterior give for a single random instrument?
\item\label{ex:iv:samesex} In the same sex design, describe the compliers. Using the accounting identity, explain why the local effect of a third child need not equal the average effect of a third child among all mothers of two.
\item\label{ex:iv:iwscale} In the structural model of Section~\ref{iv:sec:weak} take $\pi = 0$ and no controls. Show that the likelihood depends on $(\tau, \Sigma)$ only through the covariance matrix $\Omega$ of $(T, Y)$, that for each $\tau$ exactly one $\Sigma$ gives a fixed $\Omega$, and that $|\Sigma| = |\Omega|$ along the resulting curve. Show that along it, with $\Omega$ at the sample covariance matrix, the inverse Wishart prior with scale $c\,\mathrm{diag}(\Omega_{11}, \Omega_{22})$ is proportional to $\exp\{-c\,\Omega_{11}^2(\tau - b)^2/(2|\Omega|)\}$, where $b = \Omega_{12}/\Omega_{11}$ is the least squares slope. In the design of Figure~\ref{fig:weakiv}, what lengths of 90\% intervals does this predict for $c = 0.01$, 0.1, 1 and 10, and why are the lengths reported in Section~\ref{iv:sec:weak} shorter?
\item\label{ex:iv:flatK} Let the first stage have $K$ instruments, a fixed error variance $\sigma_T^2$ and a flat prior on $\pi$. Show that the concentration parameter $\mu^2 = \pi^\top Z^\top Z\pi/\sigma_T^2$ then has a prior density proportional to $(\mu^2)^{K/2 - 1}$. Compare $K = 1$ with $K = 20$, and use the comparison to explain the coverage of the diffuse structural posterior with one and with twenty weak instruments in Section~\ref{iv:sec:weak}.
\end{exercises}
